\section{A source-exact projective pitchfork for C1}
\label{sec:new-c1-projective-pitchfork}

\paragraph{Theorem and precise source domain.}
Take a bounded physical canonical C1 source with $d=1$,
$\lambda=n=0$, $m,C,A,\alpha$ positive bounded units and
$D,\beta\geq0$ infinitesimal. Assume the original actual finite
noncommuting stationary zero contact, and retain every original
source gap and intersection. Then the entire parent has a finite
sum-of-squares certificate in $K_c\otimes_{\mathbb R}K_t$ for
every independent projective test and every real or complex state.
This includes literal $D=0$ or $\beta=0$ whenever allowed by the
source domain. The two cannot vanish together at an actual contact.

\subsection{Opening-safe source entry}

Write the positive source contact as $(u_c,w_c)=(c,1/v_c)$,
and $V=w_c^2$. The retained positive-$d$ balance proof gives
$c$ a positive bounded unit and $V$ a positive infinitesimal.
Its second stationarity equation, in the notation of that proof, is
\[
 a_3=V^{-2}\{mDa_2/c^2+\beta ma_4+Da_5+\beta a_7c^2\},
 \tag{PP1}
\]
where all $a_i$ are positive bounded units. Therefore
\[
 D=V^2\delta,\quad\beta=V^2\eta,\qquad
 \delta,\eta\geq0\text{ bounded},\qquad
 \delta+\eta\text{ a positive unit}.
 \tag{PP2}
\]
No inverse of either opening is used.

Normalize the source state by $Z_c=1$. Set
\[
 p=W_c/w_c,\qquad y=Y_c/w_c,\qquad x=X_c/V.
 \tag{PP3}
\]
All three are bounded source elements. To prove this without a
shape hypothesis, restrict the reciprocal full parent PP6 to
$(W,Y,X)$ at $w=0$, and then take residues. Its positive hard form is
\[
 H_u(W,Y,X)=m(W+uY)^2+C(uW-Y)^2
       +(A+\alpha+1/2)(1+u^2)Y^2+(1+u^2)X^2/2.
 \tag{PP4}
\]
It is positive definite for every bounded real $u$. Scale the two
$W,Y$ kernel equations at the source by $w_c$, and the $X$
equation by $V$. In the unknowns $(p,y,x)$ the resulting bounded
matrix has a block triangular residue: its $p,y$ block is the
positive $W,Y$ block of PP4, and its $x$ pivot is $(1+c^2)/2$.
The terms $(\beta-D)/V=V(\eta-\delta)$ are bounded. All right
hand sides are bounded, proving PP3 by an ordinary source-field
matrix inverse with unit determinant.

Let $q_{01},q_{10},q_{21},h$ be the original Laurent coefficients.
Their first three residues are positive units because $m,C,A,\alpha$
are units. The actual zero and angular stationarity equations,
using PP2, give
\[
 \bar q_{01}=\bar q_{21}\bar c^{\,4},\qquad
 \bar h=-2\bar q_{21}\bar c^{\,2}.
 \tag{PP5}
\]
Every other nonconstant Laurent contribution apart from $q_{10}w^2$
vanishes in this residue calculation. These are equations wholly in
the coefficient field followed by its real residue map.

\subsection{A complete local normal and divided-difference argument}

It remains locally to treat testing $w$ infinitesimal or zero and
$u$ with the same positive real residue as $c$. The full reciprocal
quarter-parent is
\[
\begin{split}
 F={}&m(W-wX+uY-uwZ)^2+CD_w(uW-Y)^2+AD_uD_wY^2\\
 &+(D+1/2)(uwW+uX-wY-Z)^2
       +(\alpha+1/2)D_u(Y-wZ)^2\\
 &+(\beta+1/2)D_w(X+uZ)^2-D_uD_wZ^2/2,
 \qquad D_u=1+u^2,\quad D_w=1+w^2.
\end{split}\tag{PP6}
\]
Use three unrestricted normals and two scalar conductors
\[
 N_W=W-pwZ,\quad N_Y=Y-ywZ,\quad N_X=X-xw^2Z,
 \qquad Q=w(u-c)Z,\quad L=(w^2-V)Z.
 \tag{PP7}
\]
The centers agree with the actual source contact; no equation is
imposed on an independent state. Write the full normal expansion
\[
 F=N^TH(u,w)N+2Z\sum_iN_i\ell_i(u,w)+E(u,w^2)Z^2,
 \quad N=(N_W,N_Y,N_X)^T.
 \tag{PP8}
\]
It is an exact expansion of every term of PP6. The matrix $H$ has
bounded entries and residue PP4 at $u=\bar c$. Parity and direct
inspection of PP6 give bounded polynomials $a_W,a_Y,a_X$ such that
\[
 \ell_W=w a_W(u,b),\qquad \ell_Y=w a_Y(u,b),\qquad
 \ell_X=u(\beta-D)+b a_X(u,b),\quad b=w^2.
 \tag{PP9}
\]
Actual state stationarity gives $a_W(c,V)=a_Y(c,V)=0$ and
$\ell_X(c,w_c)=0$. Finite monic division thus gives bounded
polynomials $A_i,B_i$ with
\[
 a_i(u,b)=(u-c)A_i(u,b)+(b-V)B_i(b),\qquad i=W,Y.
\]
Consequently the entire first two normal couplings are
$Z\ell_i=A_iQ+wB_iL$.

For the third normal put
$r(u)=u(\beta-D)/V+a_X(u,V)$. This is a bounded source
polynomial by PP2, and $r(c)=0$. Hence $r=(u-c)R(u)$ for
a bounded polynomial $R$. Division in $b$ gives another bounded
polynomial $B_X$ with
\[
 Z\ell_X=\{B_X(u,b)-(u-c)R(u)\}L+wR(u)Q.
 \tag{PP10}
\]
Here the identity $V(u-c)Z=wQ-(u-c)L$ is exact. Moreover
$B_X(c,0)$ has zero residue: the source equation
$u(\beta-D)+Va_X(u,V)=0$ at $u=c$, divided by $V$,
gives $a_X(c,V)$ infinitesimal, and $B_X(c,0)$ has the same
residue. Thus the third normal has no residue coupling to $L$
or $Q$. The first two have no residue coupling to $L$.

For completeness, the scalar coefficient obeys
\[
 E(u,0)=D+\beta u^2,\qquad
 E(c,V)=E_u(c,V)=E_b(c,V)=0.
 \tag{PP11}
\]
The first assertion follows by setting $w=0$ in the entire PP6
after the graph substitution. The other three follow from the
actual full source contact and its two angle derivatives; neither
normal term is discarded in making this inference.
Define $d(u)=\delta+\eta u^2$ and the bounded polynomial
\[
 J(u,b)=\{E(u,b)-d(u)(b-V)^2\}/b.
 \tag{PP12}
\]
This notation means exact monic division: PP11--PP2 make the
constant term zero. No testing $b$ is inverted. At $(c,V)$
the value and both first derivatives of $J$ vanish. A second
finite monic Taylor division gives bounded polynomials $J_{20}$,
$J_{11},J_{02}$ with
\[
 J(u,b)=(u-c)^2J_{20}(u,b)
       +(u-c)(b-V)J_{11}(b)+(b-V)^2J_{02}(b).
 \tag{PP13}
\]
Substitution yields the complete exact scalar identity
\[
 E(u,w^2)Z^2=J_{20}(u,w^2)Q^2+wJ_{11}(w^2)QL
                     +\{d(u)+w^2J_{02}(w^2)\}L^2.
 \tag{PP14}
\]
All divisions preceded tensor evaluation, with monic divisors and
bounded source coefficients. These formulas hold literally at $w=0$.

Equations PP8--PP14 represent the entire parent by a bounded
Gram on the five rows $N_W,N_Y,N_X,Q,L$. Its $L$ pivot in
residue is $\bar\delta+\bar\eta\bar c^2>0$, and its $L$
crosses have zero residue. It remains to check the four-row hard
and angular residue; the following exact determinant computation
provides that check without taking a mixed inverse.

\subsection{Strictness of the angular residue}

At $D=\beta=0$, replace $Z$ by $T/w$ in the whole reciprocal
parent and take its polynomial value at $w=0$. This gives a real
four-state quadratic form $B_u(W,Y,X,T)$, with hard block PP4.
The original determinant identity gives
\[
 \det B_u=\frac{(1+u^2)^2}{16}
                  (q_{01}+hu^2+q_{21}u^4).
 \tag{PP15}
\]
Using PP5, its real Schur coefficient at $T$ is
\[
 g(u)=\frac{(1+u^2)^2q_{21}}{16\det H_u}
                       (u^2-c^2)^2.
 \tag{PP16}
\]
Here all entries and inverses are real residues; this calculation
proves positivity of a real matrix and does not factor a matrix in
the tensor ring. The contact centers PP3 give the unique real
kernel of $B_c$, so the four-row residue in PP8--PP14 is precisely
the quadratic Taylor coefficient in hard normals and the change
$u-c$. Its angular Schur pivot is
\[
 g''(c)/2=\frac{(1+c^2)^2q_{21}c^2}{4\det H_c}>0.
 \tag{PP17}
\]
The hard block is PP4 and is positive definite. Thus this full
four-row residue is positive definite, as is the independent $L$
pivot. The separated strict bounded-Gram lemma proves that PP6
is tensor SOS in the complete local chart. The rows may be
unbounded or linearly dependent at a literal endpoint; only the
coefficient Gram needs a strict real residue.

\subsection{Explicit testing complement and projective endpoints}

If $u,v$ are bounded original testing coordinates, use the ordinary
state matrix except near the double zero corner. At finite nonzero
real residues, the residue Laurent polynomial is
\[
 q_{21}(u^2-c^2)^2/u^2+q_{10}/v^2>0,
 \tag{PP18}
\]
so its state matrix is positive definite by the hard $Z=0$ block.
At $u=0$, $v\ne0$ the real quarter-parent is
\[
 m(vW-X)^2+(C+A+1/2)(1+v^2)Y^2
                       +\alpha(vY-Z)^2+(1+v^2)X^2/2\succ0.
\]
At $v=0$, $u\ne0$ it is
$m(X+uZ)^2+(C+1/2)(uW-Y)^2+A(1+u^2)Y^2
+\alpha(1+u^2)Z^2+(X+uZ)^2/2$, which is likewise strict.
These are residue identities; the bounded full matrices therefore
give tensor SOS on their respective charts.

When both $u,v$ are infinitesimal or zero, use the five rows
$(uW,vW,X,Y,Z)$ in the entire original finite parent. The exact
polynomial substitution is bounded and its residue is
\[
 m(vW-X)^2+(C+1/2)(uW-Y)^2+AY^2+\alpha Z^2+X^2/2\succ0
 \tag{PP19}
\]
as a coefficient Gram on those five rows. This treats that corner
with no state normalization.

For reciprocal Alice and bounded Bob, the complete core identity
RC20 applies unchanged with all $m,D,\alpha,\beta$ event packets
retained whole. Its residue RC21 is strict using only $C,A>0$.
For two reciprocal testing angles, retain those positive packets
and use RC24 for the four rows $W,X,Y,wZ$. The residue
$CW^2+AY^2+X^2/2+(Y-wZ)^2/2$ is strictly positive.
These identities allow the reciprocal coordinate to be literally zero.

For Bob reciprocal and bounded Alice, retain the whole $D,\beta$
packets, put $T=wZ$, and use the full four-row expression
\[
\begin{split}
 m(W-wX+uY-uT)^2+CD_w(uW-Y)^2
          +AD_uD_wY^2+\alpha D_u(Y-T)^2+K/2,
\end{split}
\]
where $K$ is the full core RC22 with $\rho_*=uW-Y$.
Its residue is exactly $B_u$. By PP15--PP16 it is strictly
positive whenever $u^2$ has residue different from $\bar c^2$,
including $u$ infinitesimal. The remaining matched residues are
PP7--PP17 after the exact Alice sign gauge if necessary.

Both signs of either angle are allowed by the exact simultaneous
state sign gauges of the original parent. Every normalization is
a positive square in a single source or testing field. Complex
states split into real and imaginary parts. The finite and
reciprocal homogeneous identities include all literal axes and
projective endpoints. This completes the all-tests theorem on the
stated positive-$d$ C1 projective source domain.

\paragraph{Scope of verification.}
The focused checker expands the full parent, verifies the normal
support and finite divisions with explicit contact-jet remainders,
and checks PP15 and the complete outer identities. Source
boundedness, vanishing contact jets, real-residue strictness and
the exhaustive join are the analytical arguments above. No
inference is made from a scalar certificate alone, from a
historical checklist, or from mixed multiplier cancellation.
